\sectionnonum{Išvados}
\label{sec:isvados}

\begin{enumerate}
    \item Sigmoidinis neuronas sėkmingai išmoko klasifikuoti krūties vėžio duomenis – geriausias modelis pasiekė $99{,}1\ \%$ validavimo ir $95{,}6\ \%$ testavimo tikslumą.
    \item Stochastinis gradientinis nusileidimas pasiekė mažesnę paklaidą ir aukštesnį validavimo tikslumą nei paketinis, tačiau paketinis metodas geriau generalizavo testavimo duomenims ($97{,}1\ \%$ prieš $95{,}6\ \%$).
    \item Paketinis gradientinis nusileidimas yra jautresnis mokymosi greičio parinkimui – su $\eta = 0{,}01$ validavimo tikslumas siekė tik $81{,}8\ \%$, tuo tarpu stochastinis metodas visais atvejais pasiekė $99{,}1\ \%$.
    \item Mokymo laikas abiem metodams buvo panašus, tačiau stochastinis metodas šiek tiek lėtesnis dėl dažnesnių svorių atnaujinimų.
    \item Iš $137$ testavimo įrašų $6$ buvo klasifikuoti neteisingai – $3$ nepiktybiniai navikai priskirti piktybiniams ir $3$ piktybiniai priskirti nepiktybiniams.
    \item Vienas sigmoidinis neuronas yra pakankamas pasiekti aukštą klasifikavimo tikslumą šiame duomenų rinkinyje, tačiau sudėtingesniems uždaviniams gali prireikti daugiasluoksnio neuroninio tinklo.
\end{enumerate}